\centerline{\bf \Huge Муницип ВсОШ}
\centerline{\bf \Huge 9 класс}

\begin{flushright}
        \scriptsize{}
\end{flushright}

\problem{1}
Числа $x_1, x_2, \ldots, x_n$ таковы, что $x_1 \geqslant x_2 \geqslant \ldots \geqslant x_n \geqslant 0$ и

$$
\frac{x_1}{\sqrt{1}}+\frac{x_2}{\sqrt{2}}+\ldots+\frac{x_n}{\sqrt{n}}=1
$$


Докажите, что $x_1^2+x_2^2+\ldots+x_n^2 \leqslant 1$.

\problem{2}
Вписанная в прямоугольный треугольник $A B C$ окружность касается катетов $C A$ и $C B$ в точках $P$ и $Q$ соответственно. Прямая $P Q$ пересекает прямую, проходящую через центр вписанной окружности параллельно гипотенузе, в точке $N$. Точка $M$ середина гипотенузы. Найдите величину угла $M C N$.

\problem{3}
Решите уравнение

$$
\left(x^3-7 x^2-5 x+75\right)^2+\left(x^3-8 x^2-5 x+84\right)^2=0
$$


\problem{4}
Натуральное число $n$ таково, что число $n+1$ делится на 8 . Докажите, что сумма всех делителей числа $n$, включая 1 и само число, делится на 8 .

\problem{5}
Двадцать одна девочка и двадцать один мальчик принимали участие в математическом конкурсе. Каждый участник решил не более шести задач. Для любых девочки и мальчика найдётся хотя бы одна задача, решённая обоими. Докажите, что была задача, которую решили не менее трёх девочек и не менее трёх мальчиков.